% !TeX root = ../main.tex
\section{Related Work}
\label{sec:related}

\textbf{Agentic failure attribution} has gained interest recently for ensuring LLM-based agents' reliability~\citep{qi2026trajdebug,barke2026agentrx}.
\citet{whoandwhen} formulate automated failure attribution, introduce the
\textsc{Who\&When} benchmark, and evaluate All-at-Once, Step-by-Step, and Binary Search
LLM judges; \textsc{Who\&When~Pro} scales the benchmark to 12K failed trajectories
across modalities~\citep{liu2026wwpro}, and complementary work categorizes
common multi-agent failure
modes~\citep{mast,deshpande2025trail,raj2026model,zhu2025agentdebug}.
Subsequent methods improve inference-time prompting: RAFFLES iterates a
judge--evaluator loop~\citep{raffles}, ErrorProbe pairs an analyzer with a
self-improving verifier~\citep{errorprobe}, CHIEF constructs causal dependency
graphs through multi-stage LLM reasoning~\citep{chief}, and others
restructure the context, retrieve error schemata, or reconstruct step
dependencies~\citep{echo,correct,bakish2026adaptive,falat}. Another line
trains dedicated attributors: AgenTracer and StepFinder learn from labeled or
synthetically corrupted trajectories by reinforcement learning or
temporal-semantic modeling~\citep{agentracer,zhu2026stepfinder}, and
GraphTracer learns to trace root causes over information-flow
graphs~\citep{graphtracer}. These approaches, however, rely on
\textit{computationally intensive inference or step-level error annotations},
limiting their practicality; we compare with representatives of both lines in
Table~\ref{tab:main}. OAT avoids both by one-class learning over clean
successful trajectories~\citep{oat}, but its success-only data \textit{carry
no decisive-error signal}, unlike our step-unlabeled failure trajectories,
and it underperforms \soap (Table~\ref{tab:main}).  \citet{du2024haloscope} detects hallucinations using unlabeled LLM generations but trains a classifier on top of a different projection score without the dependency propagation compared to our approach. \citet{masprism} also use
attention signals, but score them on steps with high negative log-likelihood
for two-stage localization; MASC targets online self-correction via
prototype-guided next-execution reconstruction rather than post-hoc
error attribution~\citep{shen2025masc}.

\vspace{-1em}
% Old version (pre-2026-09-08) kept below, commented out:
% \textbf{Agentic failure attribution} has gained interest recently for ensuring LLM-based agents' reliability~\citep{qi2026trajdebug,barke2026agentrx}. 
% \citet{whoandwhen} formulate automated failure attribution, introduce the
% Who\&When benchmark, and evaluate All-at-Once, Step-by-Step, and Binary Search
% LLM judges; complementary work categorizes common multi-agent failure
% modes~\citep{mast,deshpande2025trail,raj2026model}. Subsequent methods either improve inference-time prompting~\citep{echo,raffles,correct,bakish2026adaptive}
% or train dedicated attributors. For instance, AgenTracer and StepFinder instead learn
% from labeled or synthetically corrupted trajectories using reinforcement
% learning or temporal-semantic modeling~\citep{agentracer,zhu2026stepfinder}. These approaches, however, often rely on c\textit{omputationally intensive inference or large-scale step-level error annotations}, limiting their practicality for real-world deployment. A recent approach OAT mitigates this by
% learning continuous latent dynamics from clean successful
% trajectories by one-class learning~\citep{oat}, which \textit{provides no decisive error signals} in their data unlike our adopted step-unlabeled failure trajectories thus underperforms our \soap (Table~\ref{}). Several methods model cross-step dependencies~\citep{falat}. GraphTracer derives
% information-flow graphs from references to prior outputs and learns to trace
% root causes and propagation paths~\citep{graphtracer}; CHIEF construct
% causal dependency graphs during prompting through multi-stage LLM reasoning~\citep{chief}. We compare with both of them in Table~\ref{}. \cite{masprism} use attention signals for failure attribution but different from \soap, they calculate attention scores on steps with high negative log-likelihood for two-stage localization.
%
% \note{update the related work}



% \soap{} differs in its representation, dependency construction, and data
% requirements. It constructs a spectral reference space directly from unlabeled
% trajectory representations, without step labels, success-only filtering, or
% synthetic fault injection. It then derives a soft dependency graph from the
% proxy model's post-softmax attention, requiring neither explicit agent
% references nor LLM-generated causal graphs. Rather than searching the graph or
% generating a verdict, \soap{} directly propagates downstream anomaly evidence
% toward its likely upstream sources, correcting downstream contamination in a
% lightweight single-pass procedure. Experiments in
% Section~\ref{sec:experiments} compare against prompting, retrieval, trained,
% graph-based, and representation-based methods, with \soap{} outperforming the
% strongest reproduced baseline by \TODO{} without judge decoding or labeled
% reference trajectories. \textcolor{blue}{to be revised.}


% \textbf{Failure attribution in multi-agent systems.} \citet{whoandwhen} formulate automated failure attribution, release the Who\&When benchmark, and evaluate three LLM-as-judge strategies---all-at-once, step-by-step, and binary search; complementary studies catalogue why multi-agent systems fail~\citep{mast}. Two lines of work build on this. The first keeps the LLM judge but engineers the inference: ECHO organizes long traces into a hierarchical context and aggregates several role-conditioned analysts through confidence-weighted consensus voting~\citep{echo}, and CORRECT distills recurrent failures into reusable error schemata that are retrieved at inference time to guide localization~\citep{correct}. The second trains a dedicated attributor: AgenTracer constructs synthetic failure trajectories via counterfactual replay and fault injection, then fine-tunes an 8B tracer with reinforcement learning to emit the decisive step~\citep{agentracer}. All of these read the entire trajectory and return a single verdict, and the accurate ones depend on either a labeled corpus (CORRECT's schema library, AgenTracer's training set) or a strong, multi-call LLM judge. CRR differs on both axes: it scores steps from a proxy model's internal representations rather than querying a judge, and it requires no labeled trajectories at all. We also note that CORRECT observes error \emph{propagation} across agents, but uses it only as motivation; CRR turns the same observation into an explicit correction on the score field.\textcolor{blue}{please refine this by discussing with more related works.}

% \paragraph{Process reward models and step-level evaluation.} Process reward models assign a correctness score to each reasoning step and have been used to locate the earliest erroneous step in mathematical problem solving~\citep{mathshepherd,processbench}. They require dense per-step supervision, which does not exist for failure attribution---the benchmark provides only a single decisive-step label per trajectory---so a process reward model cannot be trained for this task. CRR obtains a per-step signal without any such labels.

% \paragraph{Representation-based anomaly and OOD detection.} Our base scorers transplant tools from the out-of-distribution detection literature, including norm-based scores~\citep{gradnorm} and subspace-projection scores~\citep{sal}, to the task of locating erroneous steps from a proxy model's hidden states. Our contribution is not a new anomaly score but the reframing of attribution as anomaly scoring and the causal correction layered on top.

% \paragraph{Attention as a dependency signal.} Attention weights have been used for token-level attribution, interpretability, and KV-cache compression. We instead read post-softmax attention mass as an inter-step causal-dependency estimate, converting it into row-stochastic weights that drive the residual subtraction at the heart of CRR.